\chapter{DishBrain}
\chaptermark{DishBrain}
\label{sec:dishbrain}

\section{シャーレの中のネットワークがテニスをする}

DishBrain とは、電極付きのチップ上で育てた生きたニューロンの培養が、ラケット1本の簡略化したテニスをするテストベンチに著者らがつけた名前である\cite{Kagan2022}。生きたニューロンでできたマシンの実験の中で最も議論を呼んだものであり（こうしたマシンの概観は第\ref{sec:livingmachine}章）、本書にとってはとりわけ扱いやすい。ここでは小さなネットワーク全体に手が届く。入力はすべて実験者が決め、出力はすべて記録される。これまでの章の問い、つまり信号はどう符号化されているか、誰が教えるか、プローブに何が見えるかを、目も筋肉も \nt{DOPA} ニューロンももたないネットワークに問うことができる。エンジニアが他人のテストベンチを調べるように、この実験を分解しよう。回路、入力、出力、教師、結果、論争の順である。

\section{テストベンチ}

ニューロンはマウス胎児の大脳皮質から採る（チップあたり約 $800\,000$ 個）か、ヒトの幹細胞から育てる（チップあたり約100万個）。数週間チップ上で育つうちに、ニューロンは互いに突起を伸ばし合い、自らスパイクとバーストを出し始める。その下の $8$~mm$^2$ の区画には $26\,000$ 本の白金電極が並ぶ。同時に記録できるのはそのうち最大 $1024$ 本で、刺激に使うのは実際には8本である。

培養の周りにはループが閉じている（図\ref{fig:dishloop}）。コンピュータがゲームをシミュレートする。ボールが飛び、壁で跳ね返り、ラケットが上下に動く。ボールの位置は8本の「感覚」電極の電流に変換される。二つの「運動」領域のスパイクを $10\ms$ の窓ごとに数え、ラケットを動かす。ヒットやミスのたびに、培養はもう一つの刺激、つまりフィードバックを受ける。$10\ms$ の窓はテストベンチのコンピュータのクロックであり、培養のクロックではない。ネットワーク自体は、本書のすべてのネットワークと同じく非同期で動いている。

\begin{figure}[t]
\centering
\begin{tikzpicture}[>=Stealth,
  blk/.style={draw,very thick,rounded corners=3pt,align=center,font=\small,minimum height=1.2cm},
  lab/.style={font=\scriptsize,align=center}]
\node[blk,fill=blue!6,text width=3.4cm] (C) at (0,0) {チップ上の培養\\\scriptsize $\sim10^6$ ニューロン\\\scriptsize 電極 $26\,000$ 本};
\node[blk,fill=orange!10,text width=3.0cm] (G) at (7.2,0) {コンピュータ：\\テニスゲーム};
\draw[->,very thick,blue!60!black] (G.north west) to[bend right=25] node[lab,above] {ボール：\emph{場所}は8本のどの電極か、\\\emph{頻度}は 4--40 Hz} (C.north east);
\draw[->,very thick,red!70!black] (C.south east) to[bend right=25] node[lab,below] {ラケット：「上」と「下」の領域の\\$10\ms$ ごとのスパイク数} (G.south west);
\draw[->,thick,densely dashed,green!45!black] (G.east) -- ++(0.9,0) |- ++(0,-2.6) -| node[lab,pos=0.25,below] {ヒット：予測可能な刺激、ミス：予測不能な刺激} (C.south);
\end{tikzpicture}
\caption{DishBrain のループ。青い矢印は入力で、ボールの位置は場所による符号とレート符号で書かれる（第\ref{sec:code}章）。赤は出力で、二つの領域のカウントの差がラケットを動かす。緑の破線は各ラリー後のフィードバックで、テストベンチ唯一の「教師」である。}
\label{fig:dishloop}
\end{figure}

\section{入力と出力}

\emph{入力}は第\ref{sec:code}章の二つの符号で同時に書かれる。8本のうちどの電極を刺激するかがボールの高さを表す。これは第\ref{sec:sensors}章のセンサーと同じ、場所による符号である。刺激パルスの頻度がボールの遠さを表す。ボールが遠い壁にあるときは毎秒 $4$ 回、ラケット側の壁に近づくにつれて線形に増え、毎秒 $40$ 回になる。これはレート符号で、ボールのパルスの振幅は $75$~mV である。培養はどちらの符号もあらかじめ「知って」はいない。符号は実験者が決めたもので、ネットワークはそれを読むことを学ばなければならない。

\emph{出力}は第\ref{sec:debugging}章のインターフェースと同じように読む。一方の運動領域のスパイクはラケットを上へ、もう一方は下へ押す。最も単純に書けば、各窓で
\begin{empheq}[box=\keybox]{equation}
\Delta p \;=\; c\,\bigl(n_\uparrow - n_\downarrow\bigr),
\label{eq:dishreadout}
\end{empheq}
ここで $n_\uparrow$、$n_\downarrow$ は二つの領域の $10\ms$ あたりのスパイク数、$c$ はテストベンチの係数である。これは第\ref{sec:glutcomputer}章の2線式の符号である。動きの向きを担うのは信号の符号ではなく、2本の配線のどちらが大きいかである。

この出力のノイズを見よう。ある領域が合計で毎秒 $R$ 個のスパイクを出すなら、窓 $T=10\ms$ の中には平均 $RT$ 個あり、\eqref{eq:rateerror}により相対的なばらつきは $1/\sqrt{RT}$ である。$R=1000$ スパイク/秒なら各窓で約 $30$ パーセント、$R=100$ なら100パーセントを超える。ラケットは避けがたく震え、ネットワークが学習できる方法はただ一つ、カウントの差の\emph{平均}を正しい向きにずらすことである。これは第\ref{sec:debugging}章のあらゆるインターフェースを制約するのと同じ法則であり、今回はループの両側に効いている。

\section{ドーパミンのない教師}

このテストベンチで最も興味深いのは教師である。大脳皮質ニューロンの培養には \nt{DOPA} ニューロンがなく、テストベンチは「予想より良い」という信号をまったく送らない。フィードバックは別物である。ネットワークがボールを打ち返すと、8本の感覚電極すべてに同時に短い\emph{予測可能な}刺激が与えられる。$75$~mV のパルスを毎秒 $100$ 回、$0.1\,\text{s}$ の間である。ミスすると、著者らが\emph{予測不能}と呼ぶ刺激が4秒続く。$150$~mV のパルスを毎秒 $5$ 回、ランダムな時刻に、ランダムに選んだ電極へ与える。その後、刺激のない休止が4秒あり、ボールが再スタートする\cite{Kagan2022}。

著者らは、ネットワークは自分の入力が予測可能になるように行動する、という仮説から出発した\cite{Friston2010}。エンジニアにとって意味は単純である。4秒間のランダムなノイズから逃れる方法は、培養にはただ一つ、ボールを打ち返すことしかない。同じ考えは、第\ref{sec:code}章のスパイクの節約と第\ref{sec:recognition}章のフレームの予測にも出てきた。

しかし、ここで本当に予測可能性が必要なのだろうか。もっと粗い仕組みもありうる。ミスの後の予測不能な刺激はネットワーク内の最近の変化を単に\emph{かき混ぜ}、ヒットの後の穏やかな刺激はそれに触れない、としよう。ネットワークの設定を一つのベクトル $\boldsymbol\psi$ で表し、設定 $\boldsymbol\psi$ でネットワークがミスする確率を $P_{\text{miss}}(\boldsymbol\psi)$ とする。揺さぶりが対称である、つまり $\boldsymbol\psi$ から $\boldsymbol\psi'$ へ移る確率がその逆と等しいなら、定常状態では各設定から出ていく流れと入ってくる流れが釣り合い、ネットワークが設定 $\boldsymbol\psi$ で過ごす時間の割合は
\begin{empheq}[box=\keybox]{equation}
\rho(\boldsymbol\psi) \;\propto\; \frac{1}{P_{\text{miss}}(\boldsymbol\psi)} .
\label{eq:dishdensity}
\end{empheq}
ミスが10分の1の設定は、10倍長く保たれる。重みをどちらに変えるべきかをネットワークに教える者はいない。良い設定がただ壊れにくいだけである。

これをモデルで確かめた。培養を二つの数に還元する。ボールが高さ $y$ に来たときラケットは点 $a\,y+b$ に来て、外れが $0.1$ 未満なら打ち返す。ミスの後は両方の数がランダムに揺さぶられ、ヒットの後はそのまま残る。$2000$ ラリーの間に、打ち返した割合は最初の200回の $0.21$ から最後の500回の $0.38$ に伸びた（40個の「培養」、図\ref{fig:dishmodel}）。\emph{毎回の}ラリーの後に揺さぶると、フィードバックに情報がなく、割合は約 $0.12$ にとどまる。揺さぶりがまったくなければ $0.19$ のままである。テストベンチの実験はこれと合っている。セッション内の有意な伸びは完全なフィードバックのときだけだった。刺激の代わりに無音にした場合も、セッションの終わりにはフィードバックなしより上手にプレイしたが、効果は小さかった\cite{Kagan2022}。ただし、この条件でもミスの後には長い休止が、ヒットの後には短い休止が続くので、結果による違いは完全には消えていないことに注意したい。

\begin{figure}[t]
\centering
\begin{tikzpicture}[x=1cm,y=5cm]
\draw[->,gray!70!black] (0,0) -- (10.6,0);
\draw[->,gray!70!black] (0,0) -- (0,0.5) node[above,font=\small,black] {打ち返した割合};
\foreach \v in {0.1,0.2,0.3,0.4}{\draw[gray!60!black] (-0.06,\v) -- (0.06,\v); \node[left,font=\scriptsize] at (0,\v) {\v};}
\foreach \x/\f/\l/\lab in {1.8/0.21/0.38/{ミスの後に\\ノイズ},5.2/0.13/0.11/{毎ラリーの後に\\ノイズ},8.6/0.19/0.19/{ノイズなし}}{
  \fill[blue!35] (\x-0.8,0) rectangle (\x-0.05,\f);
  \fill[red!60!black] (\x+0.05,0) rectangle (\x+0.8,\l);
  \node[below,font=\scriptsize,align=center] at (\x,-0.01) {\lab};
  \node[above,font=\scriptsize] at (\x-0.425,\f) {\f};
  \node[above,font=\scriptsize] at (\x+0.425,\l) {\l};}
\fill[blue!35] (7.4,0.47) rectangle (7.7,0.495); \node[right,font=\scriptsize] at (7.75,0.482) {最初の200回};
\fill[red!60!black] (7.4,0.41) rectangle (7.7,0.435); \node[right,font=\scriptsize] at (7.75,0.422) {最後の500回};
\end{tikzpicture}
\caption{「ミスでかき混ぜる」モデル（\texttt{dishbrain\_model.py}）：$2000$ ラリーの最初と最後で打ち返した割合。40個のモデル培養の平均。学習が進むのは、揺さぶりがミスの後に来てヒットの後には来ないときだけである。セッション内の有意な伸びが完全なフィードバックのときだけだった実験と同じである。}
\label{fig:dishmodel}
\end{figure}

\begin{keyblock}
\begin{proposition}[かき混ぜる教師]
\label{prop:dishscramble}
失敗がネットワークを揺さぶり、成功がそのままにしておくなら、ネットワークはひとりでに、誤りの少ない設定へと漂っていく。ある設定で過ごす時間は失敗の確率に反比例する\eqref{eq:dishdensity}。このような学習には、報酬信号も、その符号も、予測可能性そのものも要らない。成功の後と失敗の後でフィードバックが異なれば十分である。培養の振る舞いの変化は測定済みである。シャーレの中で働いている仕組みが入力の予測なのか、かき混ぜなのかは確定していない。
\end{proposition}
\end{keyblock}

第\ref{sec:dofa}章と比べてほしい。そこでもノイズは探索の役をしたが、ドーパミンには符号があった。誤差 $\delta$ が重みをどちらにずらすかを告げ、一歩は勾配\eqref{eq:perturbation}に沿って進んだ。ここには符号がない。あるのは「残す」か「揺さぶる」かだけである。この探索はより粗く遅いが、ネットワークへの要求は少ない。\nt{DOPA} ニューロンも、批評家も、数秒にわたる適格度トレースも要らない。

\section{測定されたことと論争点}

各ゲームは $20$ 分続き、最初の $5$ 分と最後の $15$ 分を比較した。完全なフィードバックを受けた培養では、ラリーが長くなり、即座のミスが減った。細胞なし、ゲームなし、コンピュータの「ラケット」を使った対照では、そうしたことは起こらなかった。ヒト細胞の培養は、平均するとマウスの培養より長く打ち返し続けた。改善は統計的に確かである（マウス9個とヒト11個の培養、各群100回を超えるゲーム）が、小さい。ネットワークは良いプレーヤーにはならず、ランダムより少し良くなっただけである。確かなのはセッション内の改善であり、数日にわたるセッション間の安定した学習は著者らには得られなかった\cite{Kagan2022}。同じグループの次の研究では、プレイ中の培養の活動が\emph{臨界}状態、つまり減衰と雪崩の境界、第\ref{sec:gaba}章で見たあの縁の上の生き方に近づき、近づくほどプレイが良くなることが見つかった\cite{Habibollahi2023}。

主な論争を呼んだのは数字ではなく言葉である。論文の題は、シャーレの中のニューロンが「知性」（sentience）を示すと主張した。研究者のグループが反論した。閉ループでの振る舞いの変化はまだ知性でも感覚でもなく、結論はもっと控えめに述べるべきだ、と\cite{Balci2023}。工学的な観点からは、未解決の問いが二つあり、どちらも仕組みに関するものである。一つ目：ネットワークは学習しているのか、つまり重みが長期的に変わっているのか、それともフィードバックが現在の状態をずらしているだけなのか。二つ目：本当に予測可能性が必要なのか、それとも命題\ref{prop:dishscramble}のように、成功と失敗への応答の違いがあれば十分なのか。すべての電極に手が届くテストベンチなら、両方に答えられる。おそらくそれがこのテストベンチの最大の価値である。

\section{テストベンチの仕様：再現のために}
\label{sec:dishspec}

以下に、このテストベンチを組み直すのに必要なものをすべてまとめる。装置、ループ、入力の符号化、出力の読み出し、フィードバック、実験プロトコルである。各パラメータには出典を記す。「論文」は、値を論文\cite{Kagan2022}の本文から採ったことを示す。「選択」は、論文に記載がなく、再現するには自分で決めなければならないことを示す。既定値を提案し、それをどう確かめるかを述べる。

\subsection*{装置と培養}

\begin{center}
\footnotesize
\setlength{\tabcolsep}{4pt}
\renewcommand{\arraystretch}{1.15}
\begin{tabular}{>{\raggedright}p{4.0cm}>{\raggedright}p{6.9cm}>{\raggedright\arraybackslash}p{1.6cm}}
\toprule
パラメータ & 値 & 出典 \\
\midrule
チップ & MaxOne アレイ：$8$~mm$^2$ の区画に白金電極 $26\,000$ 本 & 論文 \\
記録 & 同時に最大 $1024$ チャネル、毎秒 $20\,000$ サンプル & 論文 \\
刺激 & 技術的には最大 $32$ 電極、実験では独立な $8$ 電極 & 論文 \\
細胞 & マウス胎児の大脳皮質。幹細胞由来のヒトニューロン（2通りの培養法）。対照として HEK293T 細胞（ニューロンではない） & 論文 \\
播種 & チップあたり約 $10^6$ 個 & 論文 \\
実験時の培養日数 & マウス：$\sim10$ 日から。ヒト：培養法により $14$--$24$ 日または $\sim80$ 日 & 論文 \\
\bottomrule
\end{tabular}
\end{center}

\subsection*{ループと時間}

ループは $10\ms$（$200$ サンプル）の窓ごとに動く。窓ごとに運動領域の電極のスパイクを数え、ゲームを更新し、次の刺激を出す。培養内のスパイクから応答刺激までの遅れは約 $5\ms$ である。プレイ中、各電極の信号は遮断周波数 $100$~Hz の2次ベッセル高域通過フィルタを通る。信号の絶対値は $1$~Hz の1次ベッセル低域通過フィルタで平滑化され、スパイクの閾値はこの平滑化した絶対値に比例する。背景の標準偏差の6倍という閾値を論文が挙げているのは、個別の活動測定についてであり、ゲームについてではない。刺激が培養の応答として数えられないよう、$75$~mV を超える大きなパルスが同時に $15$ 個より多く来たときは、すべての信号を消去する。以上はすべて論文による。

\subsection*{ボールの符号化}

\begin{center}
\footnotesize
\setlength{\tabcolsep}{4pt}
\renewcommand{\arraystretch}{1.15}
\begin{tabular}{>{\raggedright}p{3.4cm}>{\raggedright}p{7.5cm}>{\raggedright\arraybackslash}p{1.6cm}}
\toprule
パラメータ & 値 & 出典 \\
\midrule
感覚領域 & $8$ 本の電極。隣り合う電極が隣り合うボール位置を表すように配置 & 論文 \\
場所による符号 & 電極番号 $k=1,\dots,8$ がラケット中心に対するボールの位置を表す & 論文 \\
どの座標か & ボールの垂直位置。再現ではラケットに対する位置とし、$y-p\in[-\tfrac12,\tfrac12]$ で $k=\lceil 8(y-p+\tfrac12)\rceil$ & 論文。式は選択 \\
頻度による符号 & $4$ から $40$~Hz の刺激頻度がラケットまでの距離を運ぶ & 論文 \\
目盛りの向き & ボールが反対側の壁にあるとき $4$~Hz、ラケット側の壁に近づくにつれ線形に $40$~Hz まで：$f=4+36\,(1-x/L)$。$x$ はラケット側の壁までの距離、$L$ はコートの幅 & 論文 \\
パルス波形 & 短い矩形の二相性パルス。正の相が先、負の相が後 & 論文 \\
振幅 & $75$~mV。抑制された細胞を無理に発火させないように選ばれた & 論文 \\
各相の長さ & 記載なし。刺激システムの標準設定を使い、実験中は変えない & 選択 \\
\bottomrule
\end{tabular}
\end{center}

\subsection*{ラケットの読み出し}

論文では運動領域は二つあり、一方のスパイクはラケットを上へ、他方は下へ動かす。細胞密度と活動の違いを補正するため、領域の寄与には重みがつけられる\cite{Kagan2022}。正確な式は示されていない。再現には規則\eqref{eq:dishreadout}、すなわち $10\ms$ の窓ごとに $\Delta p=c\,(n_\uparrow-n_\downarrow)$ を使い、安静時の平均活動で領域をそろえる重みをつける（選択）。$n_\uparrow$ と $n_\downarrow$ を、ゲーム前に測ったその領域の窓あたり平均カウントで割るのである。係数 $c$ は、平均カウント1個分の差がラケットを1窓あたりコートの高さの $1\%$ だけ動かすように選ぶ（選択）。すると一方の領域が一定に優勢なら、ラケットは $1$~s でコートを横断する。

チップ上の領域の配置は論文本文に座標で与えられておらず、模式図があるだけである。再現には重ならない電極群が三つあれば足りる。中央の感覚電極8本と、その両側の記録電極群で、一方が「上」、他方が「下」である（選択）。

\subsection*{ゲーム}

コートの大きさ、ラケットの長さ、ボールの速さは論文に記載がない。ボールが一定の速さで飛び、壁で跳ね返るとあるだけである。再現には（選択）：コート $1\times1$、右の壁に長さ $0.2$ のラケット（モデル~\texttt{models/dishbrain\_model.py} と同じく $|y-p|<0.1$ でヒット）、ボールは $1$~s でコートを横切る。ラリーで一度も打ち返さずにミスしたものを「エース」、3回を超えて連続で打ち返したラリーを「長いラリー」と数える（論文）。

\subsection*{フィードバック}

\begin{center}
\footnotesize
\setlength{\tabcolsep}{4pt}
\renewcommand{\arraystretch}{1.15}
\begin{tabular}{>{\raggedright}p{2.6cm}>{\raggedright}p{8.3cm}>{\raggedright\arraybackslash}p{1.6cm}}
\toprule
イベント & 与えるもの & 出典 \\
\midrule
ヒット & 予測可能な刺激：感覚電極 $8$ 本すべてに同時に、$75$~mV、$100$~Hz、$100\ms$。その間、通常のボールの刺激は中断する & 論文 \\
ミス & 予測不能な刺激：$150$~mV、$5$~Hz、$4$~s、ランダムな時刻に、$8$ 本からランダムに選んだ電極へ。その後 $4$~s 刺激なし、ボールはランダムな方向にスタート & 論文 \\
ランダムさの法則 & 記載なし。再現では、パルスの時刻を平均頻度 $5$~Hz のポアソン過程とする & 選択 \\
\bottomrule
\end{tabular}
\end{center}

\subsection*{プロトコルと対照}

1セッションは $20$~min で、最初の $5$~min と最後の $15$~min を比べる（論文）。結果の指標は三つ。平均ラリー長、エースの割合、長いラリーの割合である。実験条件（論文）：
\begin{itemize}
\item \emph{刺激}：上に述べた完全なループ。
\item \emph{無音}：ヒットやミスの刺激の代わりに同じ時間まったく刺激せず、その後ボールはランダムな方向にスタートする。
\item \emph{フィードバックなし}：ミスの後、ボールは単に跳ね返って飛び続け、ボール位置の刺激は続く。
\item \emph{安静}：ゲームは進み、ラケットは活動に従って動くが、刺激はまったくない。
\item \emph{対照}：培地だけのチップ。HEK293T 細胞。「コンピュータの中の培養」、つまり細胞の代わりのシミュレーション。
\end{itemize}
主系列の標本数：マウス培養 $9$ 個（$101$ セッション）、ヒト培養 $11$ 個（$138$）、培地のみのチップ $6$ 枚（$80$）、安静の培養 $20$ 個（$42$）、シミュレーション $3$ 個（$38$）、計 $399$ セッション。フィードバックの系列では、1日 $3$ セッションを $3$ 日で $486$ セッション、条件は「刺激」「無音」「フィードバックなし」（$n=27$、$17$、$15$）。最初の $5$ 分から最後の $15$ 分への平均ラリー長の伸びは、生きた培養でだけ得られた。フィードバックの系列で時間とともに有意に伸びたのは「刺激」条件だけだった。セッションの終わりには「無音」条件も「フィードバックなし」を上回ったが、効果は小さく、3日間でセッション間の安定した学習はなかった\cite{Kagan2022}。

\subsection*{テストベンチの1サイクル}

$10\ms$ ごとに、テストベンチのコンピュータは次のことを行う。
\begin{enumerate}
\item 直前の窓における二つの運動領域のスパイク数 $n_\uparrow$、$n_\downarrow$ を読み、安静時のその領域の平均カウントで正規化する。
\item ラケットを $\Delta p=c\,(n_\uparrow-n_\downarrow)$ だけ動かし、コートの端で制限する。
\item ボールを1ステップ進め、上、下、左の壁で跳ね返す。
\item ボールが右の壁に達したら、$|y-p|<0.1$ ならヒットで、ラリーのカウントを1増やし、ヒットの刺激（$100\ms$）を与える。そうでなければミスで、ラリーを記録し、ミスの刺激（$4$~s）を与え、$4$~s 休止してボールを再スタートする。
\item そうでなければ、ボールの垂直位置から電極 $k$ を、ラケット側の壁までの距離から頻度 $f$ を選び、次の窓まで電極 $k$ に $75$~mV の二相性パルスを頻度 $f$ で与える。
\end{enumerate}
これだけあれば、公表されたものと互換なテストベンチを組み、本章の中心的な問いを確かめられる。培養を教えるのは予測可能性なのか、それともヒットとミスへの応答の単なる違いなのか。そのためには論文の三つの条件に、論文にない四つ目を加えるとよい。ミスの後に、予測不能な刺激と同じ長さとエネルギーの\emph{予測可能な}刺激を与える条件である。それでも培養が学習するなら、効いているのは予測可能性ではなく違いである。


\section{まとめ}

\begin{table}[t]
\centering
\footnotesize
\setlength{\tabcolsep}{4pt}
\renewcommand{\arraystretch}{1.15}
\begin{tabular}{>{\raggedright}p{2.1cm}>{\raggedright}p{4.2cm}>{\raggedright}p{1.9cm}>{\raggedright\arraybackslash}p{4.6cm}}
\toprule
テストベンチの部位 & 仕組み & 本書の章 & わかっていること \\
\midrule
入力 & 場所（8本のどの電極か）と頻度（4--40 Hz） & \ref{sec:code}, \ref{sec:sensors} & 実験者が決めたもの \\
出力 & 二つの領域の $10\ms$ ごとのカウントの差\eqref{eq:dishreadout} & \ref{sec:glutcomputer}, \ref{sec:debugging} & 実験者が決めたもの。ノイズは\eqref{eq:rateerror}に従う \\
教師 & ヒットの後に予測可能な刺激、ミスの後に予測不能な刺激。ドーパミンはない & \ref{sec:dofa} & セッション内の有意な伸びは完全なフィードバックのときだけ。無音では効果が小さい。セッション間では不安定 \\
仕組み & 入力の予測か、ミスでのかき混ぜ\eqref{eq:dishdensity}か & \ref{sec:dofa}, \ref{sec:recognition} & 未確定。どちらも結果を説明する \\
ネットワークの状態 & プレイ中に臨界状態へ近づく & \ref{sec:gaba} & プレイの質との関係は測定済み \\
「知性」 & 著者らの主張 & --- & 示されていない。反論がある \\
\bottomrule
\end{tabular}
\caption{エンジニアの目で見た DishBrain。}
\label{tab:dishbrain}
\end{table}

DishBrain は、生きたニューロンでできたマシンの最も単純な形である。入力は符号で与えられ、出力はカウントで読まれ、教師は成功と失敗への応答の違いである（表\ref{tab:dishbrain}）。目も筋肉もドーパミンもないネットワークがほんの少しプレイを覚えた。それだけで、このネットワークは本書のアイデアを試すテストベンチになる。仕組みが予測であれ、かき混ぜであれ、第三の何かであれ、答えは第\ref{sec:debugging}章が勧めるやり方で見つかるだろう。ようやく全体が見えるマシンの上で、プローブとテスト信号と部品の切り離しを使って。

\section{問題}

\begin{exercise}[\lvlA]
\label{ex:22:1}
\emph{カウントをどれだけずらせばよいか。}二つの領域は合わせて $R_\uparrow+R_\downarrow=2000$ スパイク/秒を出し、流れはポアソン過程である。(a)~$R=1000$ と $R=100$ のとき、一つの領域の $10\ms$ あたりのカウントの相対的なばらつきはいくらか。(b)~ボールが飛ぶ $1$~s の間、変位\eqref{eq:dishreadout}は足し合わされる。平均の変位がばらつきの2倍になる最小の $R_\uparrow-R_\downarrow$ はいくらか。
\end{exercise}

\begin{exercise}[\lvlA]
\label{ex:22:2}
\emph{学習を見るには何ラリー必要か。}ラリーは独立で、打ち返した割合は二項分布に従う。打ち返した割合の伸びが差の標準偏差の2倍を超えるには、二つの期間それぞれに何ラリー必要か。(a)~$0.20$ から $0.25$ へ。(b)~モデルのように $0.21$ から $0.38$ へ。
\end{exercise}

\begin{exercise}[\lvlB]
\label{ex:22:3}
\emph{\eqref{eq:dishdensity}の導出。}設定 $\boldsymbol\psi$ は有限集合を動く。ミス（確率 $P(\boldsymbol\psi)>0$）のときネットワークは対称な確率 $K(\boldsymbol\psi,\boldsymbol\psi')$ で $\boldsymbol\psi'$ に移り、ヒットのときはとどまる。(a)~$\rho\propto1/P$ が定常分布であることを証明せよ。(b)~$P=0.9$ と $P=0.1$ の二つの設定のとき、ミスの割合はいくらか。
\end{exercise}

\begin{exercise}[\lvlB]
\label{ex:22:4}
\emph{モデルの吸収領域。}ラケットは $p=\min\bigl(1,\max(0,ay+b)\bigr)$ に、ボールは高さ $y\sim U(0,1)$ に来て、$|p-y|<h=0.1$ ならヒット。(a)~$|b|<h$ かつ $|a-1+b|<h$ ならネットワークがすべてのボールを打ち返すことを証明し、この領域の面積を求めよ。(b)~$a\sim U(-1,1)$、$b\sim U(0,1)$ のときの平均の打ち返し割合を数値的に求めよ。
\end{exercise}

\begin{exercise}[\lvlB]
\label{ex:22:5}
\emph{シミュレーション：設定の周りの柵。}\texttt{dishbrain\_model.py} のコピーで、$200$ 個の培養にそれぞれ $20\,000$ ラリーを与えよ。最後の $500$ 回で打ち返した割合はいくらか。揺さぶりのたびに設定を長方形 $a\in[-1,2]$、$b\in[-1,1]$ に切り詰めたらどうなるか。違いを問題\ref{ex:22:3}で説明せよ。
\end{exercise}

\begin{exercise}[\lvlB]
\label{ex:22:6}
\emph{入力符号の設計。}高さの誤差が $h=0.1$ 未満ならボールは打ち返される。ネットワークは入力を $T=0.25$~s で読み、頻度 $r$ までの範囲で約 $\sqrt{rT}$ 段階を区別でき、刺激は毎秒 $40$ パルス以下である。(a)~8本の電極での場所による符号で足りるか。(b)~1本の電極で高さを頻度として伝えられるか。
\end{exercise}

\begin{exercise}[\lvlC]
\label{ex:22:7}
\emph{研究：予測か、かき混ぜか。}モデルを $d$ 個の調整可能な数（$p=\sum_{i<d}a_iy^i$）に一般化せよ。打ち返し割合 $0.5$ に達するまでのラリー数は、かき混ぜと誤差の符号を使う規則\eqref{eq:perturbation}とで、$d$ とともにどう増えるか。テストベンチでのどんな実験が二つの仮説を区別するか。
\end{exercise}
